
\section{动态规划}
\subsection{多重背包(单调队列优化)}
\begin{lstlisting}[caption={}]
int main(){
    ios::sync_with_stdio(0),cin.tie(0),cout.tie(0);
    ll n,V;
    cin >> n >> V;
    vector<ll> f(V + 1),g(V + 1);
    for(int i = 1;i <= n;i++){
        ll c,w,v;
        cin >> v >> w >> c;
        g = f;
        vector<deque<ll>> q(w);
        for(int j = 0;j <= V;j++){
            if(!q[j % w].empty() && (j - q[j % w].front()) / w > c)
                q[j % w].pop_front();
            
            if(!q[j % w].empty()) f[j] = max(f[j],g[q[j % w].front()] + (j - q[j % w].front()) / w * v);
            while(!q[j % w].empty() && g[j] >= g[q[j % w].back()] + (j - q[j % w].back()) / w * v)
                q[j % w].pop_back();
            q[j % w].push_back(j);
        }
    }
    cout << *max_element(f.begin(),f.end());
    return 0;
}
\end{lstlisting}

\subsection{分组背包}
有 $n$ 组物品，一个容量为 $W$ 的背包，每组物品有若干，同一组的物品最多选一个，第 $i$ 组第 $j$ 件物品的体积为 $w[i][j]$，价值为 $v[i][j]$，求解将哪些物品装入背包，可使物品总体积不超过背包容量，且使总价值最大。

考虑每组中的某件物品选不选，可以选的话，去下一组选下一个，否则在这组继续寻找可以选的物品，当这组遍历完后，去下一组寻找。
\begin{lstlisting}[caption={}]
#include <bits/stdc++.h>
using namespace std;
const int N = 110;
int n, W, s[N], w[N][N], v[N][N], dp[N];
int main(){
    cin >> n >> W;
    for (int i = 1; i <= n; i++){
        scanf("%d", &s[i]);
        for (int j = 1; j <= s[i]; j++)
            scanf("%d %d", &w[i][j], &v[i][j]);
    }
    for (int i = 1; i <= n; i++)
        for (int j = W; j >= 0; j--)
            for (int k = 1; k <= s[i]; k++)
                if (j - w[i][k] >= 0)
                    dp[j] = max(dp[j], dp[j - w[i][k]] + v[i][k]);
    cout << dp[W] << "\n";
    return 0;
}
\end{lstlisting}

\subsection{有依赖的背包}
有 $n$ 个物品和一个容量为 $W$ 的背包，物品之间有依赖关系，且之间的依赖关系组成一颗树的形状，
如果选择一个物品，则必须选择它的 父节点，第 $i$ 件物品的体积是 $w[i]$，价值为 $v[i]$,
依赖的父节点的编号为 $p[i]$，若 $p[i]$ 等于 -1,则为 根节点。求将哪些物品装入背包中,
使总体积不超过总容量，且总价值最大。

定义 $f[i][j]$ 为以第 $i$ 个节点为根，容量为 $j$ 的背包的最大价值。那么结果就是 $f[root][W]$，为了知道根节点的最大价值，得通过其子节点来更新。所以采用递归的方式。
对于每一个点，先将这个节点装入背包，然后找到剩余容量可以实现的最大价值，最后更新父节点的最大价值即可。
\begin{lstlisting}[caption={}]
#include <bits/stdc++.h>
using namespace std;
const int N = 110;
int n, W, s[N], w[N][N], v[N][N], dp[N];
int main(){
    cin >> n >> W;
    for (int i = 1; i <= n; i++){
        scanf("%d", &s[i]);
        for (int j = 1; j <= s[i]; j++)
            scanf("%d %d", &w[i][j], &v[i][j]);
    }
    for (int i = 1; i <= n; i++)
        for (int j = W; j >= 0; j--)
            for (int k = 1; k <= s[i]; k++)
                if (j - w[i][k] >= 0)
                    dp[j] = max(dp[j], dp[j - w[i][k]] + v[i][k]);
    cout << dp[W] << "\n";
    return 0;
}
\end{lstlisting}

\subsection{数位DP}
\begin{lstlisting}[caption={}]
/* pos 表示当前枚举到第几位,sum 表示 d 出现的次数,limit 为 1 表示枚举的数字有限制,
zero 为 1 表示有前导 0,d 表示要计算出现次数的数 */
static constexpr int N = 15;
i64 dp[N][N],num[N];
i64 dfs(int pos, i64 sum, int limit, int zero, int d) {
    if (pos == 0) return sum;
    if (!limit && !zero && dp[pos][sum] != -1) return dp[pos][sum];
    i64 ans = 0;
    int up = (limit ? num[pos] : 9);
    for (int i = 0; i <= up; i++) {
        ans += dfs(pos - 1, sum + ((!zero || i) && (i == d)), limit && (i == num[pos]),
                   zero && (i == 0), d);
    }
    if (!limit && !zero) dp[pos][sum] = ans;
    return ans;
}
i64 solve(i64 x, int d) {
	std::fill(dp.begin(),dp.end(),-1);
    int len = 0;
    while (x) {
        num[++len] = x % 10;
        x /= 10;
    }
    return dfs(len, 0, 1, 1, d);
}
\end{lstlisting}


\subsection{SOSDP}
$O(N 2^N)$求解子集和问题
\begin{lstlisting}[caption={}]
for(int i = 0;i < (1 << N);i++)
	dp[i] = a[i];
for(int i = 0;i < N;i++)
	for(int mask = (1 << N) - 1;mask >= 0;mask--)
	    if(mask & (1 << i))
			dp[mask] += dp[mask ^ (1 << i)];
\end{lstlisting}


\subsection{悬线法}
\begin{lstlisting}[caption={}]
int c[N][N],u[N][N],r[N][N],l[N][N];
int main(){
	ios::sync_with_stdio(0),cin.tie(0),cout.tie(0);
	int n,m,ans = 0;
	cin >> n >> m;
	for(int i = 1;i <= n;i++)
		for(int j = 1;j <= m;j++)
			cin >> c[i][j];

	for(int i = 1;i <= n;i++){
		for(int j = 1;j <= m;j++){
			if(c[i][j]) u[i][j] = u[i - 1][j] + 1;
			if(c[i][j]) l[i][j] = l[i][j - 1] + 1;
		}
	}

	for(int i = 1;i <= n;i++)
		for(int j = m;j >= 1;j--)
			if(c[i][j]) r[i][j] = r[i][j + 1] + 1;
	
	for(int i = 1;i <= n;i++)
		for(int j = m;j >= 1;j--)

	for(int i = 1;i <= n;i++){
		for(int j = 1;j <= m;j++){
			if(c[i][j]){
				if(c[i - 1][j]){
					r[i][j] = min(r[i][j],r[i - 1][j]);
					l[i][j] = min(l[i][j],l[i - 1][j]);
				}
			}
			ans = max(ans,(r[i][j] + l[i][j] - 1) * u[i][j]);//最大矩形
			//ans = max(ans,min(r[i][j] + l[i][j] - 1,u[i][j]) * min(r[i][j] + l[i][j] - 1,u[i][j])); //最大正方形
		}
	}
	cout << ans << "\n";
	return 0;
}
\end{lstlisting}

\subsection{编辑距离}
\begin{lstlisting}[caption={}]
ll dp[2][2010];
/*
替换/跳过	删除
   插入		当前
*/
int main(){
	ios::sync_with_stdio(0),cin.tie(0),cout.tie(0);
	string a,b;
	cin >> a >> b;
	for(int i = 1;i <= b.size();i++)
		dp[0][i] = i;

	for(int i = 1;i <= a.size();i++){			
		dp[1][0] = i;
		dp[0][0] = i - 1;
		for(int j = 1;j <= b.size();j++){
			if(a[i - 1] == b[j - 1])
				dp[1][j] = dp[0][j - 1];
			else
				dp[1][j] = min(min(dp[0][j],dp[0][j - 1]),dp[1][j - 1]) + 1;
		}
		
		for(int j = 1;j <= b.size();j++){
			dp[0][j] = dp[1][j];
		}
	}

	cout << dp[1][b.size()] << "\n";
	return 0;
}
\end{lstlisting}

\subsection{最小斯坦纳树}
给定一个带权无向图 $G=(V,E)$ 和终端点集 $S \subseteq V$，求包含 $S$ 中所有点的最小权连通子图。
\begin{lstlisting}[caption={}]
const int N = 5e2 + 10,M = 5e4 + 10,mod = 1e6 + 7;
vector<PII> edge[N];
ll key[12],dp[N][M],n,m,k,vis[N];
priority_queue<PII,vector<PII>,greater<PII>> q;
void dijkstra(ll s){
	for(int i = 0;i <= n;i++)
		vis[i] = 0;	
	while(!q.empty()){
		PII t = q.top();
		q.pop();
		if(vis[t.second]) continue;
		vis[t.second] = 1;
		for(auto i : edge[t.second]){
			ll v = i.first,w = i.second;
			if(dp[v][s] > dp[t.second][s] + w){
				dp[v][s] = dp[t.second][s] + w;
				q.push({dp[v][s],v});
			}
		}
	}
}
int main(){
	ios::sync_with_stdio(0),cin.tie(0),cout.tie(0);
    cin >> n >> m >> k;
	for(int i = 1;i <= m;i++){
		ll u,v,w;
		cin >> u >> v >> w;
		edge[u].push_back({v,w});
		edge[v].push_back({u,w});
	}
	for(int i = 0;i <= n;i++)
		for(int j = 0;j <= (1LL << k);j++)
			dp[i][j] = 1e15;
	for(int i = 1;i <= k;i++){
		cin >> key[i];
		dp[key[i]][1LL << (i - 1)] = 0; 
	}
	for(ll s = 1;s <= (1LL << k) - 1;s++){
		for(int i = 1;i <= n;i++){
			for(ll sub = s & (s - 1);sub > 0;sub = s & (sub - 1))
				dp[i][s] = min(dp[i][s],dp[i][sub] + dp[i][sub ^ s]); 
			if(dp[i][s] < 1e15) q.push({dp[i][s],i});
		}
		dijkstra(s);
	}
	cout << dp[key[1]][(1LL << k) - 1] << "\n";
	return 0;
}
\end{lstlisting}